\documentclass[10pt,a4paper]{amsart}
\usepackage[margin=19mm]{geometry}
\usepackage{amsmath,amssymb,mathtools}
\usepackage[scr=rsfs,cal=euler]{mathalfa}
\usepackage{xfrac}
\usepackage[unicode,hidelinks,bookmarks=false]{hyperref}
\newcommand{\ie}{i.e.\ }
\newcommand{\cf}{cf.\ }
\newcommand{\resp}{respectively\ }
\newcommand{\Subsection}{\subsection}
\providecommand{\nobreakdash}{\nobreak\hbox{-}}
\setlength{\emergencystretch}{2em}
\multlinegap0pt
\usepackage[shortlabels]{enumitem}
\usepackage{xcolor}
\usepackage{amssymb}
\usepackage{float}
\usepackage{tikz}
\newtheorem{lem}{Lemma}
\newtheorem{thm}{Theorem}
\newtheorem{demo}{Demo}
\def\F{\mathbb{F}}
\def\Int{\mathbb{Z}}
\def\MOD{\operatorname{mod}}
\def\Non{\mathbb{N}}
\def\Pos{\mathbb{P}}
\def\fnh#1{\mathscr{R}\left({#1}\right)}
\newcommand{\mex}{\operatorname{mex}}
\newcommand{\nimadd}{\overset{*}{+}}
\newcommand{\nimsum}{{\sum}^{*}}
\title{Coin Turning Games \\on Partially Ordered Sets}
\author{Masao Ishikawa, Toyokazu Ohmoto, Hiroyuki Tagawa, Yoshiki Takayama}
\date{Source: arXiv:2602.13558v1 (CC BY 4.0)}
\begin{document}
\maketitle

\noindent\emph{Source and licence.} Coin Turning Games \\on Partially Ordered Sets. Version arXiv:2602.13558v1 (CC BY 4.0), distributed by arXiv under the Creative Commons Attribution 4.0 International licence, http://creativecommons.org/licenses/by/4.0/, as recorded in the arXiv licence metadata for this exact version and shown on the abstract page https://arxiv.org/abs/2602.13558v1. Full licence notice: \texttt{licenses/arxiv-2602.13558-CC-BY-4.0.txt}.\par\smallskip

\noindent\textit{Editorial scope.} Counted unit: the article's thm th:Grundy-subsFF-ruler (source0.tex lines 1813--1839). The article's complete original proof of it is delivered with it (source0.tex lines 2036--2170, one \texttt{proof} environment of 135 lines, which the article places away from the statement). No mathematics is written, repaired or re-derived: the statement and the proof are the article's own lines, in the article's own notation, and no retained line is substituted or re-flowed.  Retained uncounted as the source-local context the proof needs: lines 1643--1654; lines 1864--1889; lines 1883--1885; lines 1982--1984; lines 1999--2001.  Nothing was omitted from any retained span, so the package carries no omission record.  No retained line contains a citation, so the wrapper carries no bibliography and no bibliography entry.  No retained line contains a figure environment, an image-inclusion command or drawing code, so no image is delivered and the package declares no asset dependency.  Editorial formatting only, in addition: an AMS \texttt{amsart} preamble that restates the article's own declarations and macros used by the retained lines, namely the environments \texttt{lem}, \texttt{thm}, \texttt{demo} on separate counters, together with the standard packages those lines need.  The retained environments print in this document as Lemma 1, Theorem 1, Demo 1 in order of appearance, whereas the article prints the same items with the numbering of its own full document.  The complete native source of the article as distributed in the arXiv e-print for this exact version is bundled unchanged as \texttt{sources/194653/source0.tex}.
\begin{lem}\label{lem:chain-ruler-recurrence}
%
Let $f:\Pos\to\Non$ which satisfies $f(1)=1$, and
%
\begin{equation}\label{eq:chain-ruler-recurrence}
f(n)=\mex\Bigl\{\overset{n-1}{\underset{k=m}{\,\,\,\nimsum}}f(k) \,|\, m=1,2,\dots,n\Bigr\}
\end{equation}
%
for $n>1$.
Then we have $f(n)=\fnh{n}$.
%
\end{lem}

\begin{lem}\label{lem:Bn(q)-recurrence}
\begin{subequations}
%
Let \(d\) be a nonnegative integer, and let \(q\) be a prime power.
%
\begin{enumerate}[label=(\arabic*),leftmargin=2em]
%
\item\label{it:Bn(q)lem1}
Let $W\in B_n(q)$.
The Grundy value of $W$ in the ruler on $B_n(q)$ depends only on its dimension $d=\dim W$.
Accordingly, we denote the Grundy value of $W$ by $g_{q}(d)$.
%
\item\label{it:Bn(q)lem2}
For $m=0,1,\dots,d$, define
\begin{equation}\label{eq:rec:sq(d,m)-general}
s_{q}(d,m):=\overset{d-1}{\underset{k=m}{\,\,\,\nimsum}}G_{q}(d-m,k-m)\,g_{q}(k).
\end{equation}
%
Then
\begin{equation}\label{eq:rec-mex}
g_{q}(d)=\mex\left\{\,s_{q}(d,m) \,\mid\, m=0,1,\dots,d\,\right\}.
\end{equation}
\end{enumerate}
%
\end{subequations}
\end{lem}

\begin{equation}
g_{q}(d)=\mex\left\{\,s_{q}(d,m) \,\mid\, m=0,1,\dots,d\,\right\}.
\end{equation}

\begin{equation}\label{eq:def:s(d,m)}
s_{q}(d,m):=\overset{d-1}{\underset{k=m}{\,\,\,\nimsum}} G(d-m,k-m) g_{q}(k)
\end{equation}

\begin{equation}\label{eq:rec:s(d,m)}
s_{q}(d,m)=s_{q}(d,m+1)\nimadd s_{q}(d-1,m)\nimadd g_{q}(d-1).
\end{equation}

\begin{thm}\label{th:Grundy-subsFF-ruler}
\begin{subequations}
%
Let \(n\) be a nonnegative integer, and let \(q\) be a prime power.
Let $W\in B_n(q)$ be a subspace of $\F_{q}^{n}$.
The Grundy value of $W$ in $B_n(q)$ depends only on the dimension $\dim W$;
hence it is equal to the Grundy value of $\F_{q}^d$ in $B_d(q)$,
which is denoted by $g_{q}(d)$.
%
\begin{enumerate}[label=(\roman*)] %%%%%%%%%%%%%%%%%%%%%%%%%%%%%%%%%%%%%%%%%%%
	\item\label{item:q-even}
	If $q$ is even, then
	\begin{equation}
	g_{q}(d)=\fnh{d+1}.
	\end{equation}
	%
	\item\label{item:q-odd}
	If $q$ is odd, then 
	\begin{equation}
	g_{q}(d)=\MOD(d,3)+1.
	\end{equation}
\end{enumerate} %%%%%%%%%%%%%%%%%%%%%%%%%%%%%%%%%%%%%%%%%%%
%
Here, $\MOD(x,m)$ denotes the remainder when \(x\) is divided by \(m\).
%
\end{subequations}
\end{thm}

\begin{demo}{Proof of Theorem~\ref{th:Grundy-subsFF-ruler}}
%
\begin{enumerate}[leftmargin=14pt,label=(\roman*)] %%%%%%%%%%%%%%%%%%%%%%%%%%%%%%%%%%%%%%%%%%%
\item\label{item:proof-even}
When $q$ is even,
Lemma~\ref{lem:Bn(q)-recurrence} 
and the identity $G_{q}(n,r)=1$ for $0\leq r\leq n$ 
imply that the Grundy function $g_{q}(d)$ satisfies the initial condition $g_{q}(0)=1$ and, 
for $d>0$, the recurrence
%
\begin{equation*}
g_{q}(d)=\mex\Bigl\{ \,
\overset{d-1}{\underset{k=m}{\,\,\,\nimsum}} \, g_{q}(k)
\, \Big| \, 
m=0,1,\dots,d
\, \Bigr\}.
\end{equation*}
%
By Lemma~\ref{lem:chain-ruler-recurrence}, we therefore conclude that $g_{q}(d)=\fnh{d+1}$.
%
\item\label{item:proof-odd}
\begin{subequations}
%
Let $q$ be odd, and let $s_{q}(d,m)$ and $g_{q}(d)$ be defined as above.
%
To prove the theorem we define a function $H:\Int\to\Non$ by $H(x):=\MOD(x,3)+1$.
Then it clearly satisfies
%
\[
H(x+3)=H(x),\qquad H(x) \nimadd H(x+1)\nimadd H(x+2)=0
\]
%
for all $x\in\Int$.
%
To complete the proof for $q$ odd case,
it suffices to prove the following claims.
\begin{enumerate}[label=(\alph*)]
%
\item\label{it:claim:sq(d,m)}
%
For nonnegative integers $d,m\in\Non$ with $0\leq m\leq d$,
we have
%
\begin{equation}\label{eq:claim:sq(d,m)}
%
s_{q}(d,m)=\begin{cases}
0,\qquad&\text{ if $d \equiv m$ ($\MOD\,3$),}\\
H(m),\qquad &\text{ if $d \not\equiv m$ ($\MOD\,3$).}
\end{cases}
\end{equation}
%
\item\label{it:claim:gq(d)}
For all $d\in\Non$,
\begin{equation}\label{eq:claim:gq(d)}
g_{q}(d)=H(d).
\end{equation}
%
\end{enumerate}
%
\end{subequations}
%
We prove the claim by induction on $d$,
with a nested induction on $d-m$.
%
\begin{enumerate}[label=\Roman*)]
%
\item\label{it:Case:d=0}
%
If $d=0$, then necessarily $m=0$ since $0\leq m\leq d$.
By definition, the sum in \eqref{eq:def:s(d,m)} is empty, and hence $s_{q}(0,0)=0$.
%
Moreover, $g_{q}(0)=\mex\{s_{q}(0,0)\}=1$.
Therefore, both claims \ref{it:claim:sq(d,m)} and \ref{it:claim:gq(d)} hold for $d=0$.
%
%
\item\label{it:Case:d>0}
%
Let $d\geq1$,
and assume that both claims \ref{it:claim:sq(d,m)} and \ref{it:claim:gq(d)} 
hold for all smaller values of $d$, that is, for $d-1$.
%
\begin{enumerate}[label=\roman*)]
%
\item\label{it:Case:d-m=0}
%
If $d-m=0$, that is, $m=d$, then the sum in \eqref{eq:def:s(d,m)} is empty, hence $s(d,d)=0$.
This verifies \eqref{eq:claim:sq(d,m)} for $d-m=0$.
%
\item\label{it:Case:d-m>0}
%
Now suppose $d-m\geq1$, and assume that \eqref{eq:claim:sq(d,m)} holds for $d-m-1$.
%
Using the recurrence \eqref{eq:rec:s(d,m)} together with the induction hypothesis,
we compute $s_{q}(d,m)$ as
%
\begin{equation*}
\begin{cases}
H(m+1)\nimadd H(m)\nimadd H(m-1)=0\quad&\text{ if $d-m \equiv 0$ ($\MOD\,3$)}\\
0\nimadd 0\nimadd H(d-1)=H(m)\quad&\text{ if $d-m \equiv 1$ ($\MOD\,3$)}\\
H(m+1)\nimadd H(m)\nimadd H(d-1)=H(m)\quad&\text{ if $d-m \equiv 2$ ($\MOD\,3$)}
\end{cases}
\end{equation*}
%
This establishes \eqref{eq:claim:sq(d,m)} for $d-m$.
%
\end{enumerate}
%
Combining Cases~\ref{it:Case:d-m=0} and~\ref{it:Case:d-m>0} completes the proof of \eqref{eq:claim:sq(d,m)} for $d$.
%
\par
%
Next we prove \eqref{eq:claim:gq(d)} for $d$.
%
Let $S_{q}(d)=\{s_{q}(d,m)\,\mid\,m=0,1,\dots,d\}$.
Then \eqref{eq:rec-mex} gives $g_{q}(d)=\mex(S_{q}(d))$.
%
If $d=1$, then we have $S_{q}(d)=\{0,1\}$, and hence $\mex(S_{q}(1))=2=H(1)$, as claimed.
%
If $d\geq2$, then the above result implies that $H(d)\not\in S_{q}(d)$,
that $0\in S_{q}(d)$, 
and that $H(m)\in S_{q}(d)$ whenever $m\not\equiv d$ ($\MOD\,3$).
%
Consequently, $\mex(S_{q}(d))=H(d)$, which establishes \eqref{eq:claim:gq(d)} for $d$.
%
\end{enumerate}
%
Combining Cases~\ref{it:Case:d=0} and~\ref{it:Case:d>0},
we conclude that both \ref{it:claim:sq(d,m)} and \ref{it:claim:gq(d)} hold for all $d\in\Non$.
%
\end{enumerate} %%%%%%%%%%%%%%%%%%%%%%%%%%%%%%%%%%%%%%%%%%%
%
Hence, cases \ref{item:proof-even} and \ref{item:proof-odd} complete 
the proof of the theorem for both the even and odd cases.
%
\end{demo}

\end{document}
